\chapter{The Federation Model}
\label{ch:the-federation-model}

Chapter~\ref{ch:why-one-graph-is-never-enough} argued that one schema, owned by
one team, becomes the wrong number at some point nobody notices passing. The five
chapters since have built a single service well, and the one that just finished
left the schema carrying promises a client depends on. This chapter is about
what happens when the schema stops being one service's property and becomes six.

Federation is a way to present several independent GraphQL services as one. The
client sends a query to a router and gets back an answer; it does not know which
service produced which field, and it does not need to. Each service owns its
domain, publishes its own schema, and the router stitches them into a single
composed schema called the supergraph. When a query arrives, the router
decomposes it, fans out to the services that can answer each part, assembles the
pieces, and returns.

Three ideas carry the weight, and the rest of the chapter is the detail behind
them. The first is the entity: a type that exists in more than one subgraph,
identified by a key the router can read. The second is the set of directives
that describe where a field lives, which subgraph needs what, and who is allowed
to move. The third is composition: the act of taking several subgraph schemas
and producing one that is more than the sum of them, plus the rules that reject
schemas that could never work. Nothing below compiles, because nothing below is
a single service, but everything from
chapter~\ref{ch:the-first-cut-extracting-catalog} onward depends on it.

\section{The Supergraph}
\label{sec:ch06-the-supergraph}
\index{supergraph|(}%

Start with the picture that everything from here to
chapter~\ref{ch:the-gateway-market-and-the-future} hangs on.

\begin{graphqlsdl}
# -- Catalog subgraph --
type Product @key(fields: "id") {
  id: ID!
  title: String!
  price: Money!
}

# -- Reviews subgraph --
type Product @key(fields: "id") {
  id: ID!
  reviews: [Review!]!
}

# -- The supergraph, after composition --
type Product {
  id: ID!
  title: String!
  price: Money!
  reviews: [Review!]!
}
\end{graphqlsdl}

Two subgraphs, each defining a \code{Product} type with different fields. A
composer reads both schemas, notices they share a \code{@key}, and merges them
into one \code{Product} with four fields. A client that sends a query selecting
\code{title}, \code{price} and \code{reviews} gets an answer from two services
without knowing there were two. The router knew: it fetched the Catalog's fields
first, then used the \code{id} to ask Reviews for the rest.

\index{router!role}%
The router is the single point of contact, and the single point of contact is
the whole trick. The client does not address a subgraph directly. It does not
know which subgraph owns which field. It sends one document to one endpoint, and
the router's query planner produces a plan that says which subgraph answers
which selection, in which order, with which identifiers carried between them.
Chapter~\ref{ch:how-a-federated-query-actually-runs} is that plan in motion;
this chapter is the static structure that makes it possible.

\index{subgraph!definition}%
A subgraph is a normal GraphQL service with two extra responsibilities. It has
to publish its own schema without knowing about the others, and it has to answer
two queries the router will send: \code{\_service}, which returns the SDL the
router needs for composition, and \code{\_entities}, which receives a list of
type-and-key pairs and returns the matching objects. Chapter~\ref{ch:how-a-federated-query-actually-runs}
covers both in detail. For now, all that matters is that a subgraph declares its
own types and fields, marks which ones are shareable, and says which are
entities the router can ask for.

\index{supergraph!schema}%
The supergraph schema is what the client sees, and it is not one subgraph's
schema with additions. It is the output of composition: the union of every type,
every field and every directive from every subgraph, merged according to rules
that sometimes reject the whole arrangement. A field that appears in two
subgraphs without \code{@shareable} on both is a composition error. A field of
\code{Product} that a query can reach but no subgraph can resolve is a
composition error. The composer is the gatekeeper for the claim that the
supergraph works, and section~\ref{sec:ch06-how-composition-thinks} is its
rules.

\index{Mosaic!as federated supergraph}%
Mosaic, from chapter~\ref{ch:the-first-cut-extracting-catalog} onward, is six
subgraphs behind one router. Catalog owns products and their attributes.
Reviews owns reviews. Orders owns orders and their lines. Inventory owns stock.
Pricing owns prices and currency. Accounts owns customers and their profiles.
Each is a service with its own database, its own deployment, and its own slice
of the schema, and together they answer the same queries the monolith did. The
rest of Part~II is about how that arrangement is built, and this chapter is the
language it is described in.

\index{supergraph|)}%

\section{Entities and Ownership}
\label{sec:ch06-entities-and-ownership}
\index{entities|(}%

An entity is a type the router can ask for by name. That is the definition that
matters, and everything else follows from it.

\index{entities!key directive@\code{@key} directive}%
A type becomes an entity when a subgraph puts \code{@key} on
it~\autocite{apollo2026entities}. The
\code{fields} argument names the fields that together make up the key, which is
how the router will identify one instance when it needs to ask another subgraph
for more information.

\begin{graphqlsdl}
type Product @key(fields: "id") {
  id: ID!
  title: String!
  sku: String!
}
\end{graphqlsdl}

That \code{Product} exists in the Catalog subgraph, and the key says the router
can reach it by \code{id}. If a query selects \code{title} and
\code{Product.reviews} in the same document, the router fetches the
\code{title} and the \code{id} from Catalog, then sends the \code{id} to Reviews
with a query that asks for the reviews of that exact product.

The subgraph that defines the key is the entity's owner, but ownership here is
not a claim about the domain. It is a claim about who is the authority for the
entity's existence. Catalog owns \code{Product} because it is where a
\code{Product} is created; Reviews contributes fields to one but does not get to
say which products exist. That distinction matters the moment two subgraphs
disagree about whether a product is real, and chapter~\ref{ch:entity-resolution-done-right}
is where that disagreement is handled.

\subsection*{Extending an entity from another subgraph}

\index{entities!extending}%
A subgraph that wants to add fields to an entity another subgraph owns declares
the type again, with the same key, and adds its fields:

\begin{graphqlsdl}
# In the Reviews subgraph
type Product @key(fields: "id") {
  id: ID!
  reviews: [Review!]!
}
\end{graphqlsdl}

The key must match the owning subgraph's key, and the key fields must appear
here too, because the router has to fill them in from the representation it
already holds. \code{id} appears in both subgraphs, in both schemas, and in the
supergraph it appears once.

\index{entities!reference resolver}%
That arrangement works because the Reviews subgraph can answer this question:
given a \code{Product} identifier, give me back a \code{Product} with the
reviews attached. The function that answers it is the reference resolver, and
the router calls it through the \code{\_entities} query when it needs to bridge
fields from two subgraphs. A reference resolver receives a list of
representations, each a JSON object with \code{\_\_typename} and the key fields,
and returns the matching entities. Chapter~\ref{ch:how-a-federated-query-actually-runs}
shows the exact HTTP traffic; for the model, what matters is that a subgraph
extending an entity promises to resolve it from its key.

\subsection*{Multiple keys on one type}

\index{entities!multiple keys}%
An entity can declare more than one key, and the second one enables something
the first cannot:

\begin{graphqlsdl}
type Product @key(fields: "id") @key(fields: "sku") {
  id: ID!
  sku: String!
  title: String!
}
\end{graphqlsdl}

Now the router has two routes to the same entity. When a query starts from a
field that only hands back a \code{sku}, the router can still reach the rest of
\code{Product} without first resolving \code{id}. The downstream cost is that
every subgraph extending the entity has to be able to resolve it from either
key, which means a reference resolver for each. Chapter~\ref{ch:entity-resolution-done-right}
covers that cost and the batching strategy that makes it acceptable.

The second key does not create a second entity. It creates a second way to reach
the same one, and the router picks based on what it already has.

\subsection*{What is not an entity}

\index{entities!value types}%
A type without \code{@key} is a value type. It exists only in the subgraph that
defines it, cannot be extended by another subgraph, and the router has no way to
ask for it by name. Money, an address, a cursor, a page of results: these are
value types, and most types in a federated graph are value types. An entity is
the exception, and it is expensive specifically because it crosses boundaries.
Mark too many things as entities and you have promised the router can ask for
them, which means every subgraph that touches them has to write a reference
resolver. Chapter~\ref{ch:hard-modeling-problems} is about where the boundary
between entity and value type should be drawn, and the short version is that
most boundary types are entities and most nested types are not.

\index{entities|)}%

\section{Fields Across Boundaries}
\label{sec:ch06-fields-across-boundaries}
\index{federation!field-level directives|(}%

Four directives describe what happens to a field when the type it hangs on lives
in more than one subgraph. None of them make sense on a value type, because a
value type belongs to exactly one subgraph and has no second place for a field
to come from.

\subsection*{\code{@shareable}: a field more than one subgraph can resolve}

\index{federation!shareable@\code{@shareable}}%
Without \code{@shareable}, a given field of a given object type may appear in
exactly one subgraph~\autocite{apollo2026compositionrules}. Put \code{Product.title} in Catalog and again in Reviews,
without the directive, and composition fails. The rule exists because Federation
v1 did not have it, and the result was a schema that looked correct and
sometimes produced answers from the wrong service. \code{@shareable} makes the
intention explicit: this field can be resolved by more than one subgraph, and
the router may pick either.

\begin{graphqlsdl}
# In Catalog
type Product @key(fields: "id") {
  id: ID!
  title: String! @shareable
  price: Money!
}

# In Inventory
type Product @key(fields: "id") {
  id: ID!
  title: String! @shareable
  stock: Int!
}
\end{graphqlsdl}

Here \code{title} appears in both, and the router can fetch it from whichever it
is already talking to. The practical use is reducing round trips: if a query
selects \code{price} and \code{title} and \code{stock}, the router might fetch
\code{title} from Inventory alongside \code{stock} rather than making a third
call to Catalog for one field. But that is a query-planning decision the router
makes, not a guarantee you write into the schema, and
chapter~\ref{ch:how-a-federated-query-actually-runs} shows the plan that
resulted.

The cost is that both subgraphs now have to keep the field correct. If Catalog
updates \code{title} and Inventory does not, the router serves stale data with
no way to know it is stale. That is not a bug in the router. It is the
definition of a shareable field: two sources of truth that agree because you
keep them agreeing, not because the system enforces it.

\subsection*{\code{@external}: this field comes from somewhere else}

\index{federation!external@\code{@external}}%
\code{@external} marks a field as defined elsewhere in the supergraph. It
appears in a subgraph alongside \code{@requires} or \code{@provides}, never
alone. The field still appears on the type, and the subgraph still knows its
type, but the router does not ask this subgraph to resolve it.

\begin{graphqlsdl}
# In the Shipping subgraph
type Product @key(fields: "id") {
  id: ID!
  weight: Float! @external
  shippingCost: Money! @requires(fields: "weight")
}
\end{graphqlsdl}

Shipping does not own \code{weight}. Catalog does. But \code{shippingCost}
depends on it, so Shipping declares the field as \code{@external} so the router
knows to fetch it from Catalog first and hand it along. The word
\enquote{external} is from the subgraph's point of view: this field is part of
the type's schema, but it is resolved by somebody else.

\subsection*{\code{@requires}: I need these to compute mine}

\index{federation!requires@\code{@requires}}%
\code{@requires} names the external fields a resolver cannot function without.
The example above says \code{shippingCost} needs \code{weight}. When a query
selects \code{shippingCost}, the router fetches \code{weight} from Catalog first
and includes it in the representation sent to Shipping. The resolver receives
the \code{weight} and computes the cost, and it never had to ask for it.

The \code{fields} argument is a GraphQL selection set, not just a flat list. It
can reach into nested types:

\begin{graphqlsdl}
type Product @key(fields: "id") {
  id: ID!
  dimensions: Dimensions @external
  volumetricWeight: Float! @requires(fields: "dimensions { length width height }")
}
\end{graphqlsdl}

\index{federation!requires cost@\code{@requires} cost}%
The router fetches the whole selection before calling the resolver. That is the
performance guarantee and the cost. A resolver that needs one small field costs
one extra subgraph fetch for that field, and a resolver that needs a deep
traversal costs the same traversal. Chapter~\ref{ch:entity-resolution-done-right}
measures the difference.

\subsection*{\code{@provides}: I can give you extra fields from this entity}

\index{federation!provides@\code{@provides}}%
\code{@provides} is the inverse of \code{@requires} and the more surprising of
the two. It says that when this field returns an entity, the resolver can also
supply the named fields of that entity, saving the router a follow-up fetch.

\begin{graphqlsdl}
# In the Reviews subgraph
type Review @key(fields: "id") {
  id: ID!
  body: String!
  author: Customer! @provides(fields: "displayName")
}

type Customer @key(fields: "id") {
  id: ID!
  displayName: String! @external
}
\end{graphqlsdl}

Reviews returns a \code{Customer} as part of a \code{Review}'s author. The
router would normally fetch \code{displayName} from Accounts separately.
\code{@provides(fields: "displayName")} tells the router that the resolver
already has it, so the router can skip the call. The subgraph is making a
promise it has to keep: if it says it provides \code{displayName}, the resolver
had better include it in the returned \code{Customer}, or the client gets null
and the router did not ask anybody else.

\index{federation!provides vs shareable@\code{@provides} vs \code{@shareable}}%
This is not the same as \code{@shareable}. \code{@shareable} says a field can
live in two subgraphs. \code{@provides} says a resolver for one field can supply
another field of the entity it returns, and that other field is still owned by
one subgraph. The router receives it from Reviews and does not fetch it from
Accounts, but Accounts is still where it is defined.

The directive is an optimisation and it is the one that most often goes wrong.
Getting it right means the resolver returns more than it was asked for, which
GraphQL deliberately does not do anywhere else. Chapter~\ref{ch:entity-resolution-done-right}
has the measured numbers and the failure modes.

\index{federation!field-level directives|)}%

\section{Who Owns What}
\label{sec:ch06-who-owns-what}
\index{federation!structural directives|(}%

The previous section was about fields that live on two sides of a boundary at
once. This section is about directives that change who is on which side.

\subsection*{\code{@override}: moving a field from one subgraph to another}

\index{federation!override@\code{@override}}%
A field that lives in one subgraph can be migrated to another, and
\code{@override} is how you tell the router it has moved:

\begin{graphqlsdl}
# The first version: Catalog owns `price`
# Catalog subgraph:
type Product @key(fields: "id") {
  id: ID!
  price: Money!
}

# Pricing subgraph:
type Product @key(fields: "id") {
  id: ID!
}

# Later: Pricing takes over `price`
# Pricing subgraph, after the change:
type Product @key(fields: "id") {
  id: ID!
  price: Money! @override(from: "catalog")
}
\end{graphqlsdl}

\index{federation!override progressive migration@\code{@override} progressive migration}%
The \code{from} argument names the subgraph that used to resolve the field. The
router sees it and changes where it sends queries for \code{price}, immediately,
for every client. There is no deprecation window, no dual-resolution period, and
no rollback mechanism beyond deploying the old schema again. That is a sharp
tool, and chapter~\ref{ch:strangling-the-monolith} is about using it safely:
progressive migration, percentage-based traffic shifting, and the moment it is
worth pulling rather than the easier path of keeping the field in two places.

The migrated field has to match in name and type. Pricing cannot decide
\code{price} is now \code{Float} when Catalog declared it \code{Money}. That is
not an override; it is a different field, and composition rejects it.

\subsection*{\code{@interfaceObject}: contributing to an interface from a distance}

\index{federation!interfaceObject@\code{@interfaceObject}}%
\index{entities!entity interfaces}%
An interface can be an entity. Put \code{@key} on an interface, and every object
type implementing it is an entity too, reachable through the interface's key.
Then a subgraph that does not define the interface itself can contribute fields
to every entity implementing it, all at once, without naming each type.

\begin{graphqlsdl}
# In the Media subgraph
interface Media @key(fields: "id") {
  id: ID!
  url: String!
}

type Image implements Media @key(fields: "id") {
  id: ID!
  url: String!
  width: Int!
  height: Int!
}

# In the Search subgraph
type Image @key(fields: "id") {
  id: ID!
}

type Media @interfaceObject @key(fields: "id") {
  id: ID!
  relevanceScore: Float!
}
\end{graphqlsdl}

\index{federation!interfaceObject mechanics@\code{@interfaceObject} mechanics}%
The Search subgraph never declared \code{relevanceScore}
on \code{Image}. It declared it on \code{Media}, marked
as \code{@interfaceObject}, and the composer adds the field to every type
implementing \code{Media}. If a
seventh type implements \code{Media} next month, Search's contribution applies
to it automatically, without a schema change in Search. The subgraph still needs
a reference resolver for each implementing type, because the router still sends
representations, and the resolver has to know which type it is being asked for.

\index{federation!interfaceObject vs Federation v1@\code{@interfaceObject} vs Federation v1}%
This directive replaced a Federation v1 pattern where every implementing type
had to be listed and extended separately. That pattern did not survive a new
implementation arriving, and \code{@interfaceObject} is the fix.

\subsection*{\code{@inaccessible}: hiding things from the public API}

\index{federation!inaccessible@\code{@inaccessible}}%
Not everything a subgraph defines belongs in the supergraph. A field that exists
only so another subgraph can \code{@require} it, or an internal type that
happens to be in the schema because the framework generated it, can be marked
\code{@inaccessible}:

\begin{graphqlsdl}
type Product @key(fields: "id") {
  id: ID!
  internalCatalogId: String! @inaccessible
  title: String!
}
\end{graphqlsdl}

The client never sees \code{internalCatalogId} in the supergraph schema, and
introspection does not show it. The field still exists in the subgraph, and
other subgraphs can still reference it in \code{@requires}. What changes is
that the public API is cleaned of it.

\index{federation!inaccessible types@\code{@inaccessible} types}%
The directive works on types, enum values and arguments too, not only on fields.
A whole type marked \code{@inaccessible} disappears from the supergraph while
its subgraph goes on using it. The most common use is hiding implementation
details that the gateway does not need to know about, and
chapter~\ref{ch:identity-and-authorization} uses it to keep internal
authorisation fields out of the routeable schema.

\subsection*{\code{@tag}: labelling things for tools}

\index{federation!tag@\code{@tag}}%
\code{@tag(name: "string")} attaches a label to any schema element. It does
nothing to composition, query planning or execution. It exists for tooling: a
registry can filter types by tag, a dashboard can colour nodes by tag, a CI
check can fail pull requests that add a field without the right tag. Mosaic uses
it sparingly, because a tag that is not read by a machine is a comment that
looks like a directive, and comments are cheaper. Chapter~\ref{ch:ci-cd-for-a-federated-graph}
is where a machine starts reading them.

\index{federation!structural directives|)}%

\section{How Composition Thinks}
\label{sec:ch06-how-composition-thinks}
\index{composition|(}%

Composition is the step between writing subgraph schemas and running a router.
It takes every subgraph's schema, checks them against each other, and either
produces a single supergraph schema or refuses with an error that names the
problem and the location.

The algorithm is chapter~\ref{ch:composition}'s subject, and it runs in a CLI
tool and in the router and in CI, often with different settings. This section is
the model behind the tool: what the composer checks, why it rejects what it
does, and what the errors are actually saying.

\subsection*{The rule that drives everything else}

\index{composition!satisfiability}%
A supergraph schema has to be
satisfiable~\autocite{apollo2026compositionrules}: for every valid query a client could
write, every field in the selection set must be resolvable by at least one
subgraph. That is the one rule from which the others follow, and it is harder
than it looks because the query that breaks it might never have been written by
a human. A client that asks for \code{Product.reviews} gets a response from
Reviews, and a client that asks for \code{Product.weight} gets a response from
Catalog, and a client that asks for both in one query gets both. The question
the composer asks is: does any plausible traversal of this schema lead to a
field nobody can resolve?

If Catalog defines \code{Product.price} and Inventory does not, and a query
starts in Inventory, the router can fetch \code{Product} from Catalog for the
\code{price}. That is satisfiable. If Catalog does not define \code{price}
either, and nothing else does, the field exists in the supergraph and nobody can
resolve it. That is unsatisfiable, and composition refuses.

\subsection*{The explicit rules}

\index{composition!rules}%
The satisfiability rule is the principle, and the compiler's checks are the
machinery. The Apollo Federation specification lists them, and these are the
ones that produce the most errors during development:

\index{composition!duplicate field}%
\medskip\noindent\textbf{One field, one owner, unless shared.} A field of an
object type may appear in exactly one subgraph, or in several if every
occurrence is marked \code{@shareable}. A field that appears twice without
\code{@shareable} on both is a composition error, and the message names the type
and the conflicting subgraphs. This is the error most people hit first.

\index{composition!incompatible types}%
\medskip\noindent\textbf{Shared fields must agree on type.} Two subgraphs
declaring \code{Product.title} both as \code{String!} is fine. One declaring
\code{String!} and the other \code{String} is not. The types must be identical,
nullability and list wrappers included. The composition algorithm compares the
fully resolved type, not the SDL text, so \code{[String]!} and
\code{[String!]!} are different.

\index{composition!external reachability@\code{@external} reachability}%
\medskip\noindent\textbf{\code{@external} fields must be reachable.} A field
marked \code{@external} in one subgraph must be resolvable by another in every
query that could reach it. If the field is in a \code{@requires} selection, the
composer checks that the route from the entity to that field exists.

\index{composition!override validity@\code{@override} validity}%
\medskip\noindent\textbf{\code{@override} must target a valid owner.} The
\code{from} argument must name a subgraph that actually defined the field before
the override. A typo in \code{from} is a composition error, and so is overriding
a field that was never defined.

\index{composition!key matching@\code{@key} matching}%
\medskip\noindent\textbf{Keys must match across subgraphs.} When Catalog defines
\code{Product @key(fields: "id")} and Reviews also defines
\code{Product @key(fields: "id")}, the keys must name the same fields of
compatible types. A key of \code{fields: "sku"} in one and \code{fields: "id"}
in the other means the entity has two separate keys, which is valid; a key of
\code{fields: "id"} of type \code{ID!} in one and \code{String!} in the other is
not.

\index{composition!interfaceObject rules@\code{@interfaceObject} rules}%
\medskip\noindent\textbf{\code{@interfaceObject} must resolve for every
implementing type.} The subgraph that applies \code{@interfaceObject} must have
a reference resolver for every concrete type implementing the interface, or the
composer cannot guarantee that a query reaching the interface will reach a
resolvable entity.

\subsection*{What composition does not check}

\index{composition!limitations}%
A list of what is outside the composer's scope is as useful as the list of what
is inside it, because everything not on it is a runtime surprise.

\begin{itemize}
\item \textbf{Resolver behaviour.} Composition checks the schema. It cannot know
  whether a resolver returns correct data, crashes, or is slow. The supergraph
  is satisfiable on paper and still broken in production, and the only
  detector is testing.

\item \textbf{Data consistency.} A shareable field is a promise two subgraphs
  agree on its value. Composition validates the types, not the values. If
  Catalog has \code{title: "Table"} and Inventory has \code{title: "Desk"},
  composition passes and the client gets whichever one the router fetched.
  Chapter~\ref{ch:strangling-the-monolith} is where this stops being
  theoretical.

\item \textbf{Performance.} A satisfiable query plan can still be slow.
  \code{@requires} that pulls a deep tree, a chain of entities six subgraphs
  long, a reference resolver without a DataLoader: the composer signs off
  on all of them. Chapter~\ref{ch:resilience-and-performance} has the
  profiling.

\item \textbf{Cycles.} Two subgraphs can each \code{@require} a field from the
  other, and composition accepts the arrangement. At runtime the router detects
  the cycle and returns an error, which is a worse position than having the
  composer say no before deployment. Chapter~\ref{ch:hard-modeling-problems}
  covers how to spot them before the schema is committed.
\end{itemize}

\subsection*{Reading a composition error}

\index{composition!error messages}%
The first time composition fails, the message looks like a stack trace from a
language you did not write. Here is a real one, paraphrased:

\begin{minted}{text}
SATISFIABILITY_ERROR: The following supergraph API query:
  { product(sku: "MOS-FRN-0001") { reviews { author { displayName } } } }
cannot be satisfied by the subgraphs because:
  - Customer.displayName is not resolvable from subgraph "reviews"
  - Customer is not defined in subgraph "reviews"
\end{minted}

The error names a query, not just a type and a field. That is the key to reading
it: the composer found a traversal path that reaches a field nobody owns, and it
printed the path. The fix is either to mark the field \code{@shareable} where it
is, to remove it from where it should not be, or to add \code{@external} and a
\code{@requires} that tells the router where to get it.

Chapter~\ref{ch:composition} and appendix~\ref{app:composition-errors} decode
each error class with the fix and the example. For now, the point is that a
composition error is answering a question you want answered before deployment:
does this schema work? The answer is sometimes no, and getting it from a
compiler is better than getting it from a client at runtime.

\index{composition|)}%

\section*{Your turn}

This chapter has no tag and no code to run. It is the floor plan for the next
six chapters, and the exercises are about reading federation SDL rather than
writing it. Open a text editor and the Apollo Federation specification at
\url{https://specs.apollo.dev/federation/v2.15/}~\autocite{apollo2026federationspec};
everything below is answerable from the directives in this chapter and the
composition rules in
section~\ref{sec:ch06-how-composition-thinks}.

\begin{enumerate}
  \item \textbf{Trace one entity across three subgraphs.} Here is a
  \code{Product} as it might look in Mosaic once three subgraphs exist:

  \begin{graphqlsdl}
# Catalog
type Product @key(fields: "id") @key(fields: "sku") {
  id: ID!
  sku: String!
  title: String!
}

# Inventory
type Product @key(fields: "id") {
  id: ID!
  stock: Int!
  warehouseLocation: String
}

# Pricing
type Product @key(fields: "id") {
  id: ID!
  price: Money!
}
  \end{graphqlsdl}

  A client sends this query:

  \begin{minted}[fontsize=\footnotesize]{graphql}
  { product(sku: "MOS-FRN-0001") { title stock price } }
  \end{minted}

  Write down which subgraph is called first, what key is used, and the
  sequence of calls the router makes. Then say what changes if Pricing declares
  \code{price} with \code{@shareable}.

  \item \textbf{Find the unsatisfiable field.} Here is a schema that will not
  compose:

  \begin{graphqlsdl}
# Catalog
type Product @key(fields: "id") {
  id: ID!
  title: String!
}

# Reviews
type Product @key(fields: "id") {
  id: ID!
  reviews: [Review!]!
  averageRating: Float! @requires(fields: "reviews { rating }")
}
  \end{graphqlsdl}

  Name the specific thing the composer rejects and why. Then write the
  \code{@external} declaration that fixes it.

  \item \textbf{Spot the migration that fails.} In the \code{@override} example
  of section~\ref{sec:ch06-who-owns-what}, Pricing takes over \code{price} with
  \code{@override(from: "catalog")}. Say what happens if the \code{from}
  argument reads \code{"inventory"} instead. Then say what happens if Pricing
  declares the field as \code{Float!} when Catalog has it as \code{Money!}.

  \item \textbf{Design the key for a shared type.} Mosaic has an
  \code{Address} value type on \code{Customer} and \code{Order}. If
  \code{Address} needed to become an entity so a Shipping subgraph could add
  fields to it, what would the key be? There is no single right answer, but say
  why the key you chose is or is not enough across three subgraphs.

  \item \textbf{Count the queries in a chain.} A client asks for:

  \begin{minted}[fontsize=\footnotesize]{graphql}
  { me { orders { items { product { title reviews { body } } } } } }
  \end{minted}

  Assume Accounts owns \code{me}, Orders owns \code{Order} and
  \code{OrderItem}, Catalog owns \code{title} on \code{Product}, and
  Reviews owns \code{body} on \code{Review}. Each owns its entity and nothing else. Count the minimum
  number of subgraph calls the router must make, and say what entity keys are
  carried between which subgraphs. Then identify the call that could be avoided
  if one subgraph \code{@provides} something.

  \item \textbf{Write a composition error in your own words.} Take the
  unsatisfiable schema from exercise~2, before you fixed it, and write the
  composition error message the way you would explain it to a teammate who has
  never seen federation. Name the query the composer would print, the field that
  cannot be resolved, and what the fix is.

  \item \textbf{Read a subgraph and find what it owes.} Open the Apollo
  Federation subgraph specification and find the section on
  \code{Query.\_entities}. Read the input and output types. Then write down, in
  one sentence, what a subgraph has to implement beyond its own schema for the
  router to be able to ask it for an entity. You are looking at
  chapter~\ref{ch:how-a-federated-query-actually-runs}'s subject.
\end{enumerate}

Exercises~1 and~5 are the ones that will be directly tested the moment
chapter~\ref{ch:the-first-cut-extracting-catalog} splits Mosaic for real. If you
can trace a query through three subgraphs on paper, you can read the query plan
the router will print a few chapters from now, and reading the plan is what
separates knowing federation from trusting it.

